\documentclass[11pt,a4paper]{article}
\usepackage[T1]{fontenc}
\usepackage[utf8]{inputenc}
\usepackage[french]{babel}
\usepackage{lmodern}
\usepackage{geometry}
\usepackage{hyperref}
\usepackage{enumitem}
\geometry{margin=2.4cm}

\title{Presentation Technique\\MediPlan: Fonctionnement et Plan de Modification Page par Page}
\author{Equipe Projet}
\date{\today}

\begin{document}

\maketitle

\section*{1) Vision rapide}
MediPlan est une application web de gestion des rendez-vous medicaux orientee reduction du no-show.

Le projet est compose de:
\begin{itemize}[leftmargin=*]
  \item un frontend Next.js (interface utilisateur),
  \item un backend FastAPI (API metier),
  \item une couche de donnees mock cote frontend pour la demo actuelle.
\end{itemize}

\section*{2) Architecture actuelle}

\subsection*{2.1 Frontend}
\begin{itemize}[leftmargin=*]
  \item Point d'entree global: \texttt{app/layout.tsx}
  \item Page de connexion: \texttt{app/page.tsx}
  \item Layout protege (app authentifiee): \texttt{app/(app)/layout.tsx}
  \item Navigation et shell: \texttt{components/layout/app-shell.tsx}
  \item Donnees de demo: \texttt{lib/mock-api.ts} + \texttt{lib/mock-data.ts}
  \item Session/role: \texttt{contexts/app-context.tsx}
\end{itemize}

\subsection*{2.2 Backend}
\begin{itemize}[leftmargin=*]
  \item API FastAPI: \texttt{backend/app/main.py}
  \item Endpoints OFS:
  \begin{itemize}
    \item \texttt{POST /appointments}
    \item \texttt{PATCH /appointments/\{id\}/cancel}
    \item \texttt{POST /waitlist/\{siteId\}/enqueue}
    \item \texttt{GET /dashboard/kpi}
    \item \texttt{POST /risk-score/recompute/\{appointmentId\}}
  \end{itemize}
\end{itemize}

\section*{3) Flux fonctionnel (comment ca marche)}
\begin{enumerate}[leftmargin=*]
  \item L'utilisateur arrive sur la page de connexion (\texttt{app/page.tsx}).
  \item Le contexte applicatif valide le login demo via \texttt{mockUsers} (\texttt{contexts/app-context.tsx}).
  \item Une fois authentifie, l'utilisateur est redirige vers \texttt{/dashboard}.
  \item Le layout \texttt{app/(app)/layout.tsx} protege les routes (si non authentifie, retour \texttt{/}).
  \item Chaque page appelle aujourd'hui \texttt{lib/mock-api.ts} pour charger/modifier les donnees.
  \item La navigation est pilotee par le role dans \texttt{components/layout/app-shell.tsx}.
\end{enumerate}

\section*{4) Pages et responsabilites}
\begin{enumerate}[leftmargin=*]
  \item \texttt{app/page.tsx}: connexion + quick login demo
  \item \texttt{app/(app)/dashboard/page.tsx}: vue operationnelle globale
  \item \texttt{app/(app)/appointments/page.tsx}: operations rendez-vous
  \item \texttt{app/(app)/agenda/page.tsx}: planning visuel jour/semaine
  \item \texttt{app/(app)/waitlist/page.tsx}: reaffectation des creneaux
  \item \texttt{app/(app)/risk-center/page.tsx}: traitement des risques no-show
  \item \texttt{app/(app)/kpi/page.tsx}: analyse KPI et tendances
  \item \texttt{app/(app)/audit/page.tsx}: tracabilite des actions
  \item \texttt{app/(app)/settings/page.tsx}: parametrage des regles/metriques
\end{enumerate}

\section*{5) Plan de modification page par page}
Objectif global: migrer progressivement des mocks frontend vers le backend FastAPI, sans casser l'experience UI.

\subsection*{Etape A - Fondations (avant les pages)}
\begin{enumerate}[leftmargin=*]
  \item Creer une couche client API unique (\texttt{lib/api-client.ts}).
  \item Ajouter une variable d'environnement \texttt{NEXT\_PUBLIC\_API\_BASE\_URL}.
  \item Definir un mode fallback mock pour la demo (si backend indisponible).
\end{enumerate}

\subsection*{Etape B - Modification page par page (ordre recommande)}

\textbf{1. Page connexion (\texttt{app/page.tsx})}
\begin{itemize}[leftmargin=*]
  \item Etat actuel: auth 100\% mock.
  \item A faire:
  \begin{itemize}
    \item garder login demo mais preparer branchement auth reelle ensuite,
    \item gerer redirection si deja authentifie.
  \end{itemize}
\end{itemize}

\textbf{2. Dashboard (\texttt{app/(app)/dashboard/page.tsx})}
\begin{itemize}[leftmargin=*]
  \item A faire:
  \begin{itemize}
    \item brancher KPI et alertes sur \texttt{GET /dashboard/kpi},
    \item conserver chart et cartes existantes.
  \end{itemize}
\end{itemize}

\textbf{3. Rendez-vous (\texttt{app/(app)/appointments/page.tsx})}
\begin{itemize}[leftmargin=*]
  \item A faire:
  \begin{itemize}
    \item brancher creation/annulation/confirmation via API backend,
    \item conserver filtres et tableau UI.
  \end{itemize}
\end{itemize}

\textbf{4. Agenda (\texttt{app/(app)/agenda/page.tsx})}
\begin{itemize}[leftmargin=*]
  \item A faire:
  \begin{itemize}
    \item brancher la vue planning sur les slots et appointments backend,
    \item garder le mode jour/semaine et les filtres existants.
  \end{itemize}
\end{itemize}

\textbf{5. Waitlist (\texttt{app/(app)/waitlist/page.tsx})}
\begin{itemize}[leftmargin=*]
  \item A faire:
  \begin{itemize}
    \item brancher envoi d'offre et acceptation/refus,
    \item synchroniser les compteurs de cartes.
  \end{itemize}
\end{itemize}

\textbf{6. Risk center (\texttt{app/(app)/risk-center/page.tsx})}
\begin{itemize}[leftmargin=*]
  \item A faire:
  \begin{itemize}
    \item brancher recompute no-show,
    \item brancher actions de rappels et confirmations.
  \end{itemize}
\end{itemize}

\textbf{7. KPI, Audit, Settings}
\begin{itemize}[leftmargin=*]
  \item KPI: brancher metriques backend avec alertes.
  \item Audit: brancher flux backend d'audit logs.
  \item Settings: brancher sauvegarde des regles et seuils.
\end{itemize}

\section*{6) Premiere modification deja demarree}
La page de connexion a ete amelioree pour rediriger automatiquement vers \texttt{/dashboard} si l'utilisateur est deja authentifie.

Fichier modifie:
\begin{itemize}[leftmargin=*]
  \item \texttt{app/page.tsx}
\end{itemize}

\section*{7) Criteres de validation par page}
Pour chaque page migree:
\begin{enumerate}[leftmargin=*]
  \item UI identique ou meilleure.
  \item Donnees chargees depuis backend en priorite.
  \item Fallback mock optionnel en dev.
  \item Gestion des erreurs reseau (toast + etat vide).
  \item Role et permissions respectes.
\end{enumerate}

\section*{8) Prochaine action immediate}
Commencer la page 2 (Dashboard):
\begin{enumerate}[leftmargin=*]
  \item creer un adaptateur de donnees dashboard,
  \item connecter \texttt{GET /dashboard/kpi},
  \item mapper le format API vers le format attendu par les composants KPI/graphs.
\end{enumerate}

\end{document}